\chapter{Anatomy of a Predictor}\label{ch:rs:anatomy-of-a-predictor}

Two researchers report the \omterm{def:rs:anatomy-of-a-predictor:ic}{information coefficient} of the same reversal signal on the same stocks: 0.038 and 0.007. Both
computed a rank correlation between minus a past return and a future return, and both computed it correctly. The first
correlated yesterday's return with tomorrow's; the second the last five days' return with the next five days'. A \omterm{def:rs:anatomy-of-a-predictor:predictor}{predictor} is
not a formula: it is a formula, a target, a horizon, a normalisation and a \omterm{def:rs:anatomy-of-a-predictor:neutralise}{neutralisation}, and a number describing it means
nothing without all five. This chapter takes a \omterm{def:rs:anatomy-of-a-predictor:predictor}{predictor} apart into those pieces, measures it with the \omterm{def:rs:anatomy-of-a-predictor:ic}{information
coefficient} and the standard errors that the measurement needs, and ends with the \omterm{def:rs:anatomy-of-a-predictor:card}{predictor card}, the one-page record that
the rest of Part~II fills for every \omterm{def:rs:anatomy-of-a-predictor:predictor}{predictor} it studies. Its data are the book's synthetic market (chapter~5), whose
one-day reversal is known to be real.

\section{The target}

\begin{definition}[Predictor, prediction target, forecast horizon]\label{def:rs:anatomy-of-a-predictor:predictor}
A \emph{predictor}\index{predictor} is a number computed for each security at a decision time from data known at that time,
intended to be correlated across securities with a later quantity, its \emph{prediction target}\index{prediction target}
(usually a return). The \emph{forecast horizon}\index{forecast horizon} is the span of the target: the return over the next
day, the next five days, the next month.
\end{definition}

\begin{omfigure}
\begin{tikzpicture}[font=\footnotesize,x=1.1cm,y=1cm]
\draw[->,>=Stealth] (0,0) -- (11.2,0) node[right] {trading days};
\foreach \x in {0,1,...,10} \draw (\x,0.08) -- (\x,-0.08);
\node[below right] at (5.05,-0.05) {$t$};
\fill[omDef!20] (0,0.25) rectangle (4.9,0.75);
\node[omDef] at (2.45,0.5) {predictor's window: data to $t$};
\fill[omThm!20] (5,-0.75) rectangle (10,-1.25);
\node[omThm] at (7.5,-1.0) {target: return over $t+1, \dots, t+h$};
\draw[very thick,black] (5,1.0) -- (5,-1.5);
\node[above] at (5,1.0) {decision at the close of $t$};
\end{tikzpicture}

\omcaption{The timing of a \omterm{def:rs:anatomy-of-a-predictor:predictor}{predictor}. Everything it uses is known at the close of the decision day $t$; its target begins
the next day. A \omterm{def:rs:anatomy-of-a-predictor:predictor}{predictor} whose window reaches past $t$, or a target that starts at $t$, has look-ahead (chapter~3).}\label{fig:rs:anatomy-of-a-predictor:timing}
\end{omfigure}

The horizon is part of the \omterm{def:rs:anatomy-of-a-predictor:predictor}{predictor}. On the synthetic market, reversal measured four ways gives four different numbers:
minus the last day's return has a mean rank IC of 0.0385 with the next day's return ($t = 17.8$), and 0.0162 with the next
five days' return; minus the last five days' return has 0.0187 with the next day and 0.0074 with the next five ($t = 2.4$).
The simulator plants a reversal of one day, and every measurement that dilutes that day, on either side, sees less of it. A
real \omterm{def:rs:anatomy-of-a-predictor:predictor}{predictor}'s horizon is unknown and must be found (chapter~13), but the lesson carries: the pair (\omterm{def:rs:anatomy-of-a-predictor:predictor}{predictor} window,
target horizon) is what is measured.

\begin{definition}[Residual return]\label{def:rs:anatomy-of-a-predictor:residual}
The \emph{residual return}\index{residual return} of a security over a period is its return minus the part explained by
chosen factor exposures, estimated on each date by a cross-sectional regression of the returns on the exposures: what is
left after the market, the industries or the styles.
\end{definition}

Three targets are common: the raw return, the return in excess of a market (or benchmark), and a \omterm{def:rs:anatomy-of-a-predictor:residual}{residual return}. For a rank
IC across securities on one date the first two are the same, since subtracting one number from every return leaves their
ranks unchanged: on the synthetic market both give 0.0074 for the five-day reversal. The residual is different, and the
next section says when it matters.

\section{Normalisation}

\begin{definition}[Cross-sectional z-score, rank transform]\label{def:rs:anatomy-of-a-predictor:normalise}
The \emph{cross-sectional z-score}\index{cross-sectional z-score} of a \omterm{def:rs:anatomy-of-a-predictor:predictor}{predictor} on a date is its value minus that date's mean
across securities, divided by that date's standard deviation, usually clipped at a few units. Its \emph{rank
transform}\index{rank transform} replaces each value by its rank on the date, scaled to lie in $(-\tfrac12, \tfrac12)$.
\end{definition}

Raw \omterm{def:rs:anatomy-of-a-predictor:predictor}{predictors} are not comparable across dates (a five-day return of 3\% is large in a calm week and small in a crash) or
across \omterm{def:rs:anatomy-of-a-predictor:predictor}{predictors} (a return and a ratio of book value to price have different units). Normalising each date's cross-section
removes both problems and turns a \omterm{def:rs:anatomy-of-a-predictor:predictor}{predictor} into a set of relative bets, which is what a market-neutral book holds. The
z-score keeps the shape of the distribution and must be clipped (or the \omterm{def:rs:anatomy-of-a-predictor:predictor}{predictor} winsorised, Book 4, chapter 15), because one
extreme name otherwise dominates the date. The \omterm{def:rs:anatomy-of-a-predictor:normalise}{rank transform} discards the shape entirely: it is robust, and the rank IC,
Spearman's correlation of Book 4, chapter 15, is the Pearson correlation of two rank-transformed variables. Nothing is free:
the \omterm{def:rs:anatomy-of-a-predictor:normalise}{rank transform} treats the difference between the first and second names like that between the 500th and the 501st.

\section{Neutralisation}

\begin{definition}[Neutralisation]\label{def:rs:anatomy-of-a-predictor:neutralise}
\emph{Neutralisation}\index{neutralisation} is the removal from a \omterm{def:rs:anatomy-of-a-predictor:predictor}{predictor}, or from a target, of its component along chosen
exposures, by taking on each date the residual of a cross-sectional regression on them: industry neutralisation removes
industry means, beta neutralisation the component along market beta.
\end{definition}

\begin{proposition}[What neutralising the target does to the IC]\label{prop:rs:anatomy-of-a-predictor:neutral}
On one date, write the target as $y = f + u$, where $f$ is its projection on the exposures and $u$ the residual, uncorrelated
with $f$. If the \omterm{def:rs:anatomy-of-a-predictor:predictor}{predictor} $s$ is uncorrelated with $f$, then
\[
\operatorname{corr}(s, u) = \operatorname{corr}(s, y)\,\sqrt{\frac{\Var(y)}{\Var(u)}}.
\]
\end{proposition}

\begin{proof}
$\Cov(s, y) = \Cov(s, u)$, since $s$ is uncorrelated with $f$. Divide by $\sqrt{\Var(s)\Var(y)}$ and by $\sqrt{\Var(s)\Var(u)}$
in turn.
\end{proof}

If the exposures explain a share $1 - \rho$ of the target's cross-sectional variance, measuring against the residual raises the
IC by $1/\sqrt\rho$: the \omterm{def:rs:anatomy-of-a-predictor:predictor}{predictor} was right about the part it could see, and the factor part was noise to it. In the default
synthetic market the market beta and ten industries explain only 10\% of the variance of five-day returns across stocks, and
neutralising changes the five-day reversal's IC from 0.0074 to 0.0064, within noise. In a variant whose industries move more
(30\% volatility a year instead of 8\%) they explain 52\%, and the proposition predicts a ratio of $1/\sqrt{0.476} = 1.45$: the IC
goes from 0.0043 to 0.0067 against the residual target, and to 0.0086 when the \omterm{def:rs:anatomy-of-a-predictor:predictor}{predictor} is industry-neutralised too, with the
$t$-statistic going from 0.8 to 3.5 (\cref{fig:rs:anatomy-of-a-predictor:neutral}). Neutralising the \omterm{def:rs:anatomy-of-a-predictor:predictor}{predictor} removes its own
industry component, which in this market carries no information about the future.

\begin{omfigure}
\begin{tikzpicture}
\begin{axis}[width=9cm,height=5cm,ybar,bar width=26pt,ylabel={mean rank IC},tick label style={font=\small},
  label style={font=\small},xmin=-0.6,xmax=2.6,ymin=0,ymax=0.011,xtick={0,1,2},
  xticklabels={{both raw},{residual target},{both neutral}},scaled y ticks=false,
  yticklabel style={/pgf/number format/fixed,/pgf/number format/precision=3},grid=major,grid style={gray!25},
  nodes near coords,point meta=explicit symbolic,every node near coord/.append style={font=\footnotesize},
  table/col sep=comma]
\addplot[fill=omDef!55,draw=omDef] table[x=k,y=ic,meta=t]{figdata/research/06-anatomy-of-a-predictor/neutral.csv};
\end{axis}
\end{tikzpicture}

\omcaption{Five-day reversal against five-day returns in a synthetic market whose industries move strongly (30\% a year): raw
\omterm{def:rs:anatomy-of-a-predictor:predictor}{predictor} and raw target, raw \omterm{def:rs:anatomy-of-a-predictor:predictor}{predictor} and residual target, industry-neutral \omterm{def:rs:anatomy-of-a-predictor:predictor}{predictor} and residual target. Labels are the
HAC $t$-statistics. Data: \texttt{firm.synthmkt} with \texttt{ind\_vol=0.30}, seed 1.}\label{fig:rs:anatomy-of-a-predictor:neutral}
\end{omfigure}

\section{The information coefficient and its standard error}

\begin{definition}[Information coefficient, rank information coefficient]\label{def:rs:anatomy-of-a-predictor:ic}
The \emph{information coefficient}\index{information coefficient} (IC) of a \omterm{def:rs:anatomy-of-a-predictor:predictor}{predictor} on a date is the correlation, across the
securities of that date, between the \omterm{def:rs:anatomy-of-a-predictor:predictor}{predictor} and its target. The \emph{rank information coefficient}\index{rank information
coefficient} is the same with ranks (Spearman's correlation). A \omterm{def:rs:anatomy-of-a-predictor:predictor}{predictor} is summarised by the time series of its ICs: their mean,
their standard deviation, and the ratio of the two.
\end{definition}

\begin{definition}[IC information ratio, quantile spread]\label{def:rs:anatomy-of-a-predictor:icir}
The \emph{IC information ratio}\index{IC information ratio} (ICIR) is the mean IC divided by its standard deviation over dates,
annualised by $\sqrt{252/h}$ for an $h$-day horizon sampled every $h$ days. The \emph{quantile spread}\index{quantile spread} is
the mean target of the securities in the \omterm{def:rs:anatomy-of-a-predictor:predictor}{predictor}'s top quantile minus that of the bottom quantile, date by date.
\end{definition}

Two facts about the IC's noise decide how long a \omterm{def:rs:anatomy-of-a-predictor:predictor}{predictor} must be observed. On one date with $N$ names and no predictability,
the IC has a standard deviation of about $1/\sqrt{N - 1}$ (0.032 for 1\,000 names); a mean IC of 0.01 is therefore invisible on any
single date and must be accumulated over hundreds. And when the target spans $h$ days but the IC is computed every day, the
targets of consecutive dates share $h - 1$ days, the ICs are strongly autocorrelated, and the naive standard error of their mean
is too small.

\begin{proposition}[Overlapping targets]\label{prop:rs:anatomy-of-a-predictor:overlap}
If the daily IC series of an $h$-day target has autocorrelations $\rho_k$ that decline linearly to zero at lag $h$ (as for a
persistent \omterm{def:rs:anatomy-of-a-predictor:predictor}{predictor} and non-overlapping daily returns), the variance of its mean over $n$ dates is approximately $h$ times the
naive $\Var/n$: the naive $t$-statistic overstates by up to $\sqrt h$. The Newey--West estimator with $h - 1$ lags (Book 4,
chapter 11) corrects it, and so does using every $h$-th date.
\end{proposition}

\begin{proof}
$\Var(\bar x) \approx \frac{\sigma^2}{n}\bigl(1 + 2\sum_{k=1}^{h-1}\rho_k\bigr)$ with $\rho_k = 1 - k/h$; the sum is $(h - 1)/2$, so the
bracket is $h$.
\end{proof}

For the five-day reversal against five-day targets, the naive $t$-statistic is 3.85 and the corrected one 2.39: the ratio 1.61 is
below $\sqrt5 = 2.24$ because the reversal is not persistent. Using only every fifth date gives 1.31, a valid but less efficient
estimate from a fifth of the data. For the one-day reversal against one-day returns there is no overlap: 0.0385, $t = 17.8$, an
annualised ICIR of 6.4 once the \omterm{def:rs:anatomy-of-a-predictor:predictor}{predictor} is neutralised, a number no real daily \omterm{def:rs:anatomy-of-a-predictor:predictor}{predictor} reaches after costs; the simulator
plants it strongly so that later chapters' tools have something to find.

\begin{omfigure}
\begin{tikzpicture}
\begin{axis}[width=5.8cm,height=4.8cm,ybar,bar width=7pt,xlabel={day $k$ after the decision},ylabel={rank IC with day $k$'s return},
  tick label style={font=\small},label style={font=\small},xmin=0.3,xmax=10.7,ymin=-0.01,ymax=0.045,xtick={1,2,3,4,5,6,7,8,9,10},
  scaled y ticks=false,yticklabel style={/pgf/number format/fixed,/pgf/number format/precision=2},grid=major,grid style={gray!25},
  table/col sep=comma]
\addplot[fill=omDef!55,draw=omDef] table[x=k,y=ic]{figdata/research/06-anatomy-of-a-predictor/lags.csv};
\end{axis}
\end{tikzpicture}\hspace{3mm}%
\begin{tikzpicture}
\begin{axis}[width=5.8cm,height=4.8cm,xlabel={years},ylabel={cumulative IC},tick label style={font=\small},
  label style={font=\small},xmin=0,xmax=10,ymin=0,grid=major,grid style={gray!25},table/col sep=comma]
\addplot[omDef,very thick] table[x=year,y=cum]{figdata/research/06-anatomy-of-a-predictor/cumic.csv};
\end{axis}
\end{tikzpicture}

\omcaption{The one-day reversal on the synthetic market. Left: its rank IC with the return of each of the next ten days; the
information is spent on the first day (0.0385), the second has 0.0008. Right: the cumulative sum of its daily IC over ten years,
neutralised to beta and industries; a straight line is a stable \omterm{def:rs:anatomy-of-a-predictor:predictor}{predictor}. Data: \texttt{firm.synthmkt}, seed 1.}\label{fig:rs:anatomy-of-a-predictor:lags}
\end{omfigure}

\section{The predictor card}

\begin{definition}[Predictor card]\label{def:rs:anatomy-of-a-predictor:card}
A \emph{predictor card}\index{predictor card} is the one-page record of a \omterm{def:rs:anatomy-of-a-predictor:predictor}{predictor}: its definition as a formula; its inputs and
the \omterm{def:rs:point-in-time-data-and-the-biases:bitemporal}{knowledge time} of each; its rationale; its horizon and measured half-life; its normalisation and \omterm{def:rs:anatomy-of-a-predictor:neutralise}{neutralisation}; its known
failure modes; its sources. The measured statistics come from a tested script, never from memory.
\end{definition}

The card is the unit in which Part~II of this book, and a research group, accumulates knowledge: a \omterm{def:rs:anatomy-of-a-predictor:predictor}{predictor} without a card has
not been studied, only tried. Its fields force the questions of this chapter to be answered before the \omterm{def:rs:anatomy-of-a-predictor:predictor}{predictor} is used. Here is
the first.

\begin{predictorcard}{One-day reversal}\label{pred:rs:anatomy-of-a-predictor:reversal}
\sfield{Definition} $s_{i,t} = -r_{i,t}$, the day's return with the sign changed, ranked across the universe and neutralised to
market beta and industries.
\sfield{Inputs and timestamps} The closes of days $t-1$ and $t$, known at the close of $t$; betas and industries as of $t$.
\sfield{Rationale} Liquidity providers are paid, through a reversal, for absorbing demand for immediacy (chapter~1).
\sfield{Horizon and half-life} One day. On the synthetic market the rank IC is 0.0385 with day $t+1$ and 0.0008 with day $t+2$:
the half-life is under a day.
\sfield{Normalisation} Cross-sectional rank; residual on beta and industry dummies. Neutralised, its mean rank IC is 0.040 with an
annualised ICIR of 6.4 ($t = 20.1$) on the synthetic market, which plants it.
\sfield{Failure modes} Moves on news (earnings days) do not revert; the daily turnover costs more than the IC earns in liquid
large stocks; crowding by many liquidity providers.
\sfield{Sources} N.~Jegadeesh (1990); B.~N.~Lehmann (1990); \texttt{rs\_anatomy.card()} for the numbers.
\end{predictorcard}

\section{Tutorial: measuring a predictor}\label{tut:rs:anatomy-of-a-predictor:measure}

\begin{tutorial}
\textbf{Goal.} Measure reversal on the synthetic market with every choice of this chapter made explicitly, and fill its card.
\textbf{End state:} \cref{fig:rs:anatomy-of-a-predictor:neutral,fig:rs:anatomy-of-a-predictor:lags}; the horizon table; the
$t$-statistics 3.85 (naive) and 2.39 (corrected); \cref{pred:rs:anatomy-of-a-predictor:reversal}.
\begin{tutsteps}
\item \textbf{Neutralise.} A cross-sectional regression per date on a constant and the exposures; names with a missing value are
skipped, not filled.
\omcode{code/firm/predictor/firm_predictor.py}{53}{66}{Per-date residuals on the exposures.}\label{lst:rs:anatomy-of-a-predictor:residual}
\item \textbf{Summarise.} The IC series' mean, its information ratio and two $t$-statistics, the second with $h - 1$ Newey--West
lags for overlapping targets.
\omcode{code/firm/predictor/firm_predictor.py}{107}{118}{The IC summary.}\label{lst:rs:anatomy-of-a-predictor:summary}
\item \textbf{Run} \texttt{horizon\_ics()}, \texttt{study()} and \texttt{card()}; then \texttt{neutralisation} on a market simulated
with \texttt{ind\_vol=0.30}, and \texttt{fig\_anatomy.py}.
\end{tutsteps}
\textbf{What to change next.} Replace the \omterm{def:rs:anatomy-of-a-predictor:normalise}{rank transform} by a z-score clipped at 3, then at 10, and watch the IC's standard deviation;
neutralise to the size exposure as well and check that the reversal survives.
\end{tutorial}

\section{Build: the predictor toolkit}\label{bld:rs:anatomy-of-a-predictor:predictor}

\begin{build}
\bfield{Purpose} One implementation of targets, normalisers, \omterm{def:rs:anatomy-of-a-predictor:neutralise}{neutralisation} and the IC for every \omterm{def:rs:anatomy-of-a-predictor:predictor}{predictor} of the book, so that
numbers from different chapters (and, later, Books 8 and 9) are comparable.
\bfield{Interface} \texttt{forward\_return(ret, h)}, \texttt{excess(target, market)}, \texttt{residualise(y, X)} and
\texttt{neutralise}, \texttt{zscore(x, clip)}, \texttt{rank\_transform(x)}, \texttt{ic\_series(signal, target, method)},
\texttt{ic\_summary(ic, h, periods)}, \texttt{quantile\_spread(signal, target, q)}, \texttt{PredictorCard}.
\bfield{Rules} Panels are (dates, names) with NaN for absent names; a \omterm{def:rs:anatomy-of-a-predictor:predictor}{predictor} at row $t$ is known at the close of $t$ and its target
starts at $t+1$; missing values are skipped per date, never imputed; a date with fewer than 20 names has no IC.
\bfield{Acceptance tests} \texttt{code/firm/predictor/tests/}: compounding and gaps in forward returns; clipping and ranks;
residuals orthogonal to the exposures; an IC of 0.05 recovered with the $1/\sqrt N$ noise; a naive $t$ more than 2.5 times the HAC $t$
for a persistent \omterm{def:rs:anatomy-of-a-predictor:predictor}{predictor} and overlapping 20-day targets; a card's fields.
\bfield{Stretch} Weighted ICs (by liquidity); ICs by group (industry, size bucket); a bootstrap of the IC's mean over dates.
\end{build}

\begin{omsources}
\item R.~C.~Grinold and R.~N.~Kahn, \emph{Active Portfolio Management}, 2nd ed., McGraw-Hill, 2000.
\item E.~E.~Qian, R.~H.~Hua and E.~H.~Sorensen, \emph{Quantitative Equity Portfolio Management}, Chapman \& Hall/CRC, 2007.
\item W.~K.~Newey and K.~D.~West, ``A simple, positive semi-definite, heteroskedasticity and autocorrelation consistent covariance
matrix'', \emph{Econometrica} 55(3), 1987.
\item N.~Jegadeesh, ``Evidence of predictable behavior of security returns'', \emph{Journal of Finance} 45(3), 1990; B.~N.~Lehmann,
``Fads, martingales, and market efficiency'', \emph{Quarterly Journal of Economics} 105(1), 1990.
\end{omsources}

\section{Exercises}

\begin{exercise}[$\star$]\label{exo:rs:anatomy-of-a-predictor:1}
A universe has 800 names. What is the standard deviation of a single date's IC for a useless \omterm{def:rs:anatomy-of-a-predictor:predictor}{predictor}, and how many independent
dates are needed for a mean IC of 0.01 to reach $t = 2$?
\end{exercise}

\begin{exercise}[$\star$]\label{exo:rs:anatomy-of-a-predictor:2}
A \omterm{def:rs:anatomy-of-a-predictor:predictor}{predictor}'s daily rank IC against the next day's raw return is 0.02. What is its rank IC against the next day's return in
excess of the equal-weighted market? And against the return in excess of each stock's own beta times the market?
\end{exercise}

\begin{exercise}[$\star$]\label{exo:rs:anatomy-of-a-predictor:3}
A daily IC series has mean 0.012 and standard deviation 0.08. What is the annualised ICIR?
\end{exercise}

\begin{exercise}[$\star\star$]\label{exo:rs:anatomy-of-a-predictor:4}
Industry factors explain 40\% of the cross-sectional variance of monthly returns, and a stock-selection \omterm{def:rs:anatomy-of-a-predictor:predictor}{predictor} with no industry
content has an IC of 0.03 against raw returns. What IC should it show against industry-residual returns?
\end{exercise}

\begin{exercise}[$\star\star$]\label{exo:rs:anatomy-of-a-predictor:5}
A 20-day target is scored every day for a very persistent \omterm{def:rs:anatomy-of-a-predictor:predictor}{predictor}. By the proposition, by how much can the naive $t$-statistic
overstate? What $t$ does a naive 6.0 correspond to?
\end{exercise}

\begin{exercise}[$\star\star$]\label{exo:rs:anatomy-of-a-predictor:6}
Give one reason to prefer a rank IC to a Pearson IC and one reason to prefer the Pearson IC, in terms of what a portfolio built from
the \omterm{def:rs:anatomy-of-a-predictor:predictor}{predictor} earns.
\end{exercise}

\begin{exercise}[$\star\star\star$]\label{exo:rs:anatomy-of-a-predictor:7}
\emph{Coding.} With \texttt{ic\_series} and \texttt{ic\_summary}, measure the one-day reversal against the next day's return
separately on days when the market's absolute return was in its top decile and on the other days. Which days carry the \omterm{def:rs:anatomy-of-a-predictor:predictor}{predictor}?
\end{exercise}

\begin{exercise}[$\star\star\star$]\label{exo:rs:anatomy-of-a-predictor:8}
\emph{Find the flaw.} ``Our 60-day \omterm{def:rs:anatomy-of-a-predictor:predictor}{predictor} has a mean IC of 0.04 against 60-day forward returns, computed daily over five years:
1\,260 dates, so $t = 0.04/(0.10/\sqrt{1\,260}) = 14$.''
\end{exercise}

\section{Problem: Two ICs for One Signal}

\begin{problem}[{Weekend problem --- what the choices of a measurement do to a predictor's numbers}]\label{pb:rs:anatomy-of-a-predictor:1}
The synthetic market with its default configuration (1\,000 names, ten years, seed 1), and a variant with industry volatility of
30\% a year. The \omterm{def:rs:anatomy-of-a-predictor:predictor}{predictor} is reversal: minus a past return.

\medskip\noindent\textbf{Part I --- Horizon.}
\begin{enumerate}
\item What is the rank IC of one-day reversal against the next day's return, and its $t$-statistic?
\item And of five-day reversal against the next five days' return?
\item What are the two mixed pairs (one-day against five, five-day against one)?
\item Which of the four is the planted effect, and why are the others smaller?
\item What is the one-day reversal's IC with the return of the second day after the decision?
\end{enumerate}

\medskip\noindent\textbf{Part II --- Targets and \omterm{def:rs:anatomy-of-a-predictor:neutralise}{neutralisation}.}
\begin{enumerate}[resume]
\item Why do the raw and the market-excess targets give the same rank IC?
\item What share of five-day return variance do beta and industries explain in the default market?
\item In the strong-industry variant, what share, and what ratio does the proposition predict?
\item What are the three ICs of \cref{fig:rs:anatomy-of-a-predictor:neutral} and their $t$-statistics?
\item Why does neutralising the \omterm{def:rs:anatomy-of-a-predictor:predictor}{predictor} add to neutralising the target?
\end{enumerate}

\medskip\noindent\textbf{Part III --- Standard errors.}
\begin{enumerate}[resume]
\item What are the naive and corrected $t$-statistics of the five-day reversal against five-day targets?
\item Why is their ratio below $\sqrt5$?
\item What does using every fifth date give, and why is it lower still?
\item What is the standard deviation of one date's IC with 1\,000 names and no predictability?
\end{enumerate}

\medskip\noindent\textbf{Part IV --- The card.}
\begin{enumerate}[resume]
\item What are the card's mean IC, annualised ICIR and $t$-statistic?
\item What does the card say about the half-life, and from which numbers?
\item Which field of the card would change first if the market's liquidity providers became more numerous?
\item State the \emph{named result}: the five-day reversal's naive and HAC $t$-statistics for overlapping targets, and the IC gained by
neutralising target and \omterm{def:rs:anatomy-of-a-predictor:predictor}{predictor} in the strong-industry market.
\item What should a research report say when it quotes an IC?
\item In one sentence: what is a \omterm{def:rs:anatomy-of-a-predictor:predictor}{predictor}?
\end{enumerate}
\end{problem}

\section{Interview questions}

\begin{interviewq}[$\star$ \iqroles{researcher}]\label{iq:rs:anatomy-of-a-predictor:1}
What is an \omterm{def:rs:anatomy-of-a-predictor:ic}{information coefficient}, and what values are good?
\end{interviewq}

\begin{interviewq}[$\star\star$ \iqroles{researcher, mle}]\label{iq:rs:anatomy-of-a-predictor:2}
You compute a 20-day \omterm{def:rs:anatomy-of-a-predictor:predictor}{predictor}'s IC every day and get $t = 8$. Do you believe it?
\end{interviewq}

\begin{interviewq}[$\star\star$ \iqroles{researcher}]\label{iq:rs:anatomy-of-a-predictor:3}
Why neutralise a \omterm{def:rs:anatomy-of-a-predictor:predictor}{predictor} to industries? When would you not?
\end{interviewq}

\begin{interviewq}[$\star\star$ \iqroles{researcher, trader}]\label{iq:rs:anatomy-of-a-predictor:4}
A \omterm{def:rs:anatomy-of-a-predictor:predictor}{predictor} has an IC of 0.05 against raw returns and 0.02 against \omterm{def:rs:anatomy-of-a-predictor:residual}{residual returns}. What is going on?
\end{interviewq}

\begin{interviewq}[$\star\star$ \iqroles{mle, researcher}]\label{iq:rs:anatomy-of-a-predictor:5}
Rank IC or Pearson IC for a machine-learned return forecast? Why?
\end{interviewq}

\begin{interviewq}[$\star\star\star$ \iqroles{researcher}]\label{iq:rs:anatomy-of-a-predictor:6}
Derive the standard deviation of a single date's IC under no predictability, and use it to size the history needed to detect a
mean IC of 0.02.
\end{interviewq}
